\section{Qubit e Notazioni}

Un \textbf{Qubit} è il sistema quantistico più semplice, i cui stati possono essere descritti da uno spazio vettoriale complesso a due dimensioni. In generale, questo è dovuto al fatto che "la misurazione" starebbe nel capire in che verso sta spinnando l'elettrone in base ad una stessa "direzione". Ma questo è fuori dagli scopi del corso.

\subsection{Notazione di Dirac}
La \textbf{notazione di Dirac} (o \textbf{BraKet}) è un modo per simboleggiare i vettori che ci permette di "gestire" in maniera più coerente quello che è lo studio della computazione quantistica.
Considerando \( c_1, c_2, \dots, c_n \in \mathbb{C}\), allora si definiscono le seguenti notazioni.
Le notazioni di base sono quella della \textbf{Ket} (Vettore colonna):
\[
\ket{\varphi} =  
\begin{bmatrix}
c_1 \\
c_2 \\
\vdots \\
c_n    
\end{bmatrix}
\]

E della \textbf{Bra} (Vettore Riga), è il coniugato trasposto di \(\ket{\varphi}\):
\[
\bra{\varphi} = 
\begin{bmatrix}
c_1^* & c_2^* & \dots & c_n^*
\end{bmatrix}
\]

Date le precedenti notazioni, sì definisce il \textbf{prodotto interno} come:
\[
\bra{x}\ket{y} = \braket{x | y}
\]

Mentre si definisce il \textbf{prodotto esterno} come:
\[
\ket{x}\bra{y}
\]

\subsection{Qubit}
Nel caso del Qubit, essendo lo spazio di dimensione 2, si andrà a considerare quella che è la base canonica dello spazio, ovvero: \(\mathcal{B} = 
\left\{
\left[
\begin{smallmatrix}
1 \\
0    
\end{smallmatrix}
\right]
,
\left[
\begin{smallmatrix}
0 \\
1
\end{smallmatrix}
\right]
\right\}
\).
Tale base può essere anche scritta seguendo la seguente notazione:
\[
\ket{0} = \begin{bmatrix}1 \\ 0\end{bmatrix} 
\]
\[
\ket{1} = \begin{bmatrix}0 \\ 1\end{bmatrix}
\]

Pertanto possiamo scrivere la base del nostro spazio degli stati come:
\[
\mathcal{B} = \left\{ \ket{0}, \ket{1} \right\}
\]

Pertanto, uno stato quantistico può essere visto come la \textbf{superposizione} di stati, il cui risultato è un altro stato quantistico valido.
Quindi, uno stato \(\ket{\varphi} \in \mathbb{C}^2\) del Qubit, può essere visto come:
\[
\ket{\varphi} = \alpha_0 \ket{0} + \alpha_1 \ket{1}
\]

Dove, \(|\alpha_0|^2\) è la probabilità di ottenere esito \(\ket{0}\) dalla misura sulla base precedentemente presentata; mentre \(|\alpha_1|^2\) è la probabilità di misurare lo stato \(\ket{1}\). Dalla teoria della probabilità, quindi:
\[
\alpha_0, \alpha_1 \in \mathbb{C} : |\alpha_0|^2 + |\alpha_1|^2 = 1
\]
La somma dei moduli delle possibilità all'interno di un sistema quantistico dev'essere 1.

Altra base molto importante per la quantum computation è la \textbf{base di Hadamard}, ovvero:
\[
\mathcal{H} = \left\{ \ket{+},\ket{-}\right\}
\]

Dove sarà possibile esprimere la superposizione come:
\[
\ket{\varphi} = \beta_0 \ket{+} + \beta_1 \ket{-}
\]

Gli elementi della base di Hadamard possono essere espressi anche come elementi della base canonica, quindi è possibile andare a definire le seguenti relazioni:
\[
\ket{+} = \frac{1}{\sqrt{2}}\left(\ket{0} + \ket{1}\right)
\]
\[
\ket{-} = \frac{1}{\sqrt{2}}\left(\ket{0} - \ket{1}\right)
\]

